\documentclass[UTF8]{ctexart}
\usepackage{xeCJK}
\usepackage{geometry}
\geometry{left=2cm,right=2cm,top=2.5cm,bottom=2.5cm}
\setlength{\headheight}{12.7pt}

\usepackage{titlesec}
\titleformat{\section}
  {\normalfont\large\bfseries}
  {\thesection}
  {1em}
  {}
\titlespacing*{\section}
  {0pt}
  {3.5ex plus 1ex minus .2ex}
  {2.3ex plus .2ex}

\usepackage{amsmath}
\usepackage{amssymb}
\usepackage{amsthm}
\usepackage{enumitem}
\setlist[enumerate]{itemsep=0pt,topsep=0pt,parsep=0pt,leftmargin=0pt,itemindent=2.5em}
\setlist[itemize]{itemsep=0pt,topsep=0pt,parsep=0pt,leftmargin=0pt,itemindent=2.5em}
\usepackage{fancyhdr}
\pagestyle{fancy}
\fancyhf{}
\usepackage{graphicx}
\usepackage{hyperref}
\usepackage{indentfirst}
\usepackage{lmodern}


\begin{document}
\setlength{\abovedisplayskip}{1pt}
\setlength{\belowdisplayskip}{1pt}

\section{加法原理}

加法原理用于计算有限不交并的势.

\vspace{10pt}

\hypertarget{sec:定理1.1}{}
\noindent\textbf{定理1.1}

任给集合$A,B$, 如果$A\cap B=\varnothing$, 那么$|A\cup B|=|A|+|B|$.

\noindent{\kaishu 证明.}

记$A,B$的势分别为$\kappa,\lambda$, 则显然$
A\cup B\,\approx\,A\sqcup B\,\approx\,\kappa\sqcup\lambda\,
\approx\,\kappa+\lambda$. \hfill $\qedsymbol$

\vspace{10pt}

\hypertarget{sec:定理1.2}{}
\noindent\textbf{定理1.2}

任给有限集$X$和$X$上的映射$f$, 有$\displaystyle
\left|\bigsqcup_{x\in X}f(x)\right|=\sum_{x\in X}|f(x)|$.

\noindent{\kaishu 证明.}

不妨设$X$非空, 取双射$\varphi:|X|\to X$. 对$n<|X|$归纳证明
$\displaystyle\left|\,\bigcup_{i=0}^nf\varphi(i)\times\{\varphi(i)\}\,\right|=
\sum_{i=0}^n|f\varphi(i)|$.

首先, $n=0$时
\[
    \left|\,\bigcup_{i=0}^0f\varphi(i)\times\{\varphi(i)\}\,\right|=
    |\,f\varphi(0)\times\{\varphi(0)\}\,|=|f\varphi(0)|=
    \sum_{i=0}^0|f\varphi(i)|\,;\\[5pt]
\]
其次, 假设命题对$n$成立且$n+1<|X|$, 那么
\vspace{3pt}
\[\begin{aligned}
    \left|\,\bigcup_{i=0}^{n+1}f\varphi(i)\times\{\varphi(i)\}\,\right|&=
    \left|\,\bigcup_{i=0}^nf\varphi(i)\times\{\varphi(i)\}\,\right|+
    |\,f\varphi(n+1)\times\{\varphi(n+1)\}\,|\\[3pt]&=
    \sum_{i=0}^n|f\varphi(i)|+|f\varphi(n+1)|=
    \sum_{i=0}^{n+1}|f\varphi(i)|,\\[6pt]
\end{aligned}\]
归纳成立. 取$n=|X|-1$即得$\displaystyle\left|\bigsqcup_{x\in X}f(x)\right|=
\left|\bigcup_{i=0}^{|X|-1}f\varphi(i)\times\{\varphi(i)\}\right|=
\sum_{i=0}^{|X|-1}|f\varphi(i)|=\sum_{x\in X}|f(x)|$. \hfill $\qedsymbol$

\vspace{13pt}

\hypertarget{sec:定理1.3}{}
\noindent\textbf{定理1.3}

任给集合$X$和它的一个划分$S$, 如果$S$有限, 那么$\displaystyle|X|=\sum_{s\in S}|s|$.

\noindent{\kaishu 证明.}

对任意的$\displaystyle(x,s)\in\bigsqcup_{s\in S}s$定义
$\varphi(x,s)=x$, 易证$\displaystyle\varphi:\bigsqcup_{s\in S}s\to X$是双射. 从而
\vspace{3pt}
\begin{equation*}
    |X|=\left|\bigsqcup_{s\in S}s\right|=\sum_{s\in S}|s|.
\tag*{\qedsymbol}\end{equation*}

\end{document}